\documentclass[10pt]{article}
\usepackage[letterpaper,margin=0.66in]{geometry}
\usepackage{microtype}
\usepackage{amsmath}
\usepackage{booktabs}
\usepackage{enumitem}
\usepackage{xcolor}
\usepackage{hyperref}
\hypersetup{colorlinks=true,linkcolor=blue!45!black,citecolor=blue!45!black,urlcolor=blue!45!black,pdftitle={The Scheduler Free-Energy Certificate},pdfauthor={Justin D. Hart}}
\setlength{\parindent}{0pt}
\setlength{\parskip}{2.5pt}
\setlist{nosep,leftmargin=*}
\raggedbottom
\emergencystretch=2em

\title{\vspace{-1.2em}\textbf{The Scheduler Free-Energy Certificate}\\
\large A one-step audit of anti-starvation research selection}
\author{Justin D. Hart\\\small Viridis Research}
\date{8 August 2026\\\small Final post-Aristotle reconciled manuscript}

\begin{document}
\maketitle

\begin{abstract}
Autonomous research programs need both productive selection and reproducible
safety rails.  We audit a deterministic anti-starvation rule in which every
unserved research area ages by one and the selected area resets to zero.  For
inverse temperature $\beta>0$, define the soft maximum
$F_\beta(a)=\beta^{-1}\log\sum_i e^{\beta a_i}$.  We show by an exact
choice-gap identity that serving a maximum-age area minimizes the next-step
$F_\beta$, uniquely when the maximum is unique.  An adversarial standard-library
Python gate passed 22 of 22 checks, including 78,100 exhaustive and 10,000
randomized cases, tie and zero-temperature-information controls, stable
large-$\beta$ evaluation, an unequal-cost counterexample, and replay against the
current Viridis state.  The locked META transition moves total staleness from
139 to 155 and makes area 13 the unique age-22 choice for Run 120.  The audit
also finds that the governing policy names an undefined scoring variable and
contains no deterministic tie-break.  We package a versioned proposal but do
not adopt it.  Separately, Aristotle returned six Lean theorems whose frozen
statements were rebuilt and audited with zero proof holes, no forbidden
constructs, and only the permitted foundational axioms.  The result is a
formally verified one-step equal-cost scheduling certificate, not a long-run
productivity theorem, policy adoption, or execution authority.
\end{abstract}

\textbf{Keywords:} research scheduling; age of information; log-sum-exp;
anti-starvation; reproducible policy; Lyapunov potential

\section{Introduction}

A nightly research engine must choose one object while preserving a much wider
portfolio.  Selection pressure rewards depth and topical demand; fairness
pressure prevents quiet areas from becoming permanently invisible.  A safety
rail can serve that purpose only if it has two properties: its local action
should improve a declared measure, and its implementation should be reproducible
from versioned state rather than inherited prose.

Freshness scheduling has a mature theory.  Age of Information (AoI) was
introduced to distinguish timely status from throughput and queueing delay
\cite{kaul2012}.  Highest-age service is optimal in a symmetric broadcast model,
while MaxWeight and Whittle policies address more general unreliable settings
\cite{kadota2018}.  Queue-state-driven MaxWeight control has a still earlier
Lyapunov lineage \cite{tassiulas1992}.  Recent work combines scheduling with
learning under unknown and nonstationary service statistics
\cite{yang2023}, and statistical curriculum learning analyzes task order through
its effect on learning risk \cite{cohen2024}.  These results prevent an inflated
claim: prioritizing an old item is not new by itself.

Our contribution is narrower and operational.  We derive an exact one-step
certificate for the severe-starvation rail used by the Viridis nightly engine,
apply it to the locked Run 119 state, and audit whether the current policy file
fully determines future selections.  This is a scheduled META run under a
proposal-only rule.  No governing policy is edited.

\section{Preregistration and scope}

The hypotheses, assumptions, falsifiers, and executable gate were fixed before
this manuscript was drafted.

\subsection{Hypotheses}

\begin{enumerate}
\item For every $\beta>0$, resetting a maximum-age component minimizes the
next-step log-sum-exp staleness potential.
\item The locked META transition raises the live checksum from 139 to 155 and
makes area 13 the unique severe-starvation choice for Run 120.
\item The version-1 policy is not a total executable selector because it does
not define \texttt{scoring\_staleness} or deterministic tie-breaking.
\item A draft total-order proposal leaves the actual Run 119/120 replay
unchanged while resolving a deliberately tied state.
\end{enumerate}

The gate would fail on any analytic, exhaustive, or randomized counterexample;
any positive $\beta$ that changes the next choice; any existing policy
definition contradicting hypothesis 3; or any actual-state replay drift.

\subsection{Assumptions}

Ages are nonnegative integers measured before a run.  A normal research slot has
equal unit cost: all ages increment, then the selected area resets to zero.  A
META run increments every age and resets none.  Scientific importance, tier,
cross-pollination demand, lens fit, uncertain return, and unequal effort costs
are excluded from the certificate.  These are not minor omissions: they are the
reason the result cannot be promoted into a general research-allocation theorem.

\section{One-step free-energy certificate}

Let $a=(a_1,\ldots,a_n)$ have nonnegative integer entries.  If area $j$ is served, define
\begin{equation}
a_i^{(j)}=\begin{cases}0,&i=j,\\a_i+1,&i\ne j.\end{cases}
\label{eq:transition}
\end{equation}
For $\beta>0$, define
\begin{equation}
Z_\beta(a)=\sum_{i=1}^n e^{\beta a_i},\qquad
F_\beta(a)=\frac{1}{\beta}\log Z_\beta(a).
\label{eq:free}
\end{equation}
The second quantity is the smooth maximum of the ages.  It satisfies
$\max_i a_i\le F_\beta(a)\le\max_i a_i+\log(n)/\beta$.

\textbf{Proposition 1 (reset identity).}
For any service choice $j$,
\begin{equation}
Z_\beta(a^{(j)})
=1+e^\beta\left[Z_\beta(a)-e^{\beta a_j}\right].
\label{eq:reset}
\end{equation}
Consequently, for choices $j$ and $k$,
\begin{equation}
Z_\beta(a^{(j)})-Z_\beta(a^{(k)})
=e^\beta\left(e^{\beta a_k}-e^{\beta a_j}\right).
\label{eq:gap}
\end{equation}

\emph{Derivation.}  The selected exponential is replaced by $e^0=1$ and
every unselected exponential is multiplied by $e^\beta$.  Subtracting two
candidate expressions cancels all common terms.

\textbf{Theorem 1 (maximum age minimizes next-step soft maximum).}
For every $\beta>0$, $F_\beta(a^{(j)})$ is minimal exactly when
$a_j=\max_i a_i$.  The minimizer is unique if and only if the maximum age is
unique.

\emph{Derivation.}  Equation~\eqref{eq:gap} has the opposite ordering of
$a_j$, and the logarithm is strictly increasing.  Equal maximum ages give equal
minimal values.  This is the entire theorem; its value is as an inspectable
certificate tied to the implemented transition.

The thermodynamic language is an analogy, not an empirical physical claim.
$\beta$ tunes sensitivity to the tail: at large $\beta$, $F_\beta$ approaches
the maximum age; at small positive $\beta$, selection differences flatten.  At
$\beta=0$, the unscaled partition function is $n$ for every choice, a negative
control showing that positivity is load-bearing.

\section{Application to the locked scheduler}

The source-state SHA-256 is
\texttt{f329b1a6d39e52c154a5b97aa9166512ed0c64d106396f9f7763ae82c3a9286e}.
It locks Run 119 as META $\times$ Thermodynamic and records a total staleness of
139.  The META transition increments sixteen counters and resets none, giving
155.  Table~\ref{tab:ages} shows the active frontier.

\begin{table}[h]
\centering
\small
\caption{Largest post-META ages and binding precedence.}
\label{tab:ages}
\begin{tabular}{lll}
\toprule
Area & Age & Status \\
\midrule
13 Computational Theory & 22 & severe-starvation threshold \\
12 AI Wavefunction Collapse & 20 & below severe threshold \\
06 Entropy-Driven Learning & 16 & Tier-2 floor \\
07 Gaian Systems & 15 & Tier-2 floor \\
\bottomrule
\end{tabular}
\end{table}

Severe starvation precedes tier floors.  Area 13 is the sole age-22 component,
so Theorem 1 and the stated rail agree for every tested $\beta$ from $10^{-6}$
to $100$.  Lens rotation after Thermodynamic gives Stewardship.  Area 13 was
last visited at Run 097; its Run-120 separation is 23, exceeding the 16-run
combination guard.  The locked next directive is therefore
\textbf{Run 120: area 13 Computational Theory $\times$ Stewardship}.

\section{Executable-policy audit}

The governing policy SHA-256 is
\texttt{fdee0f1dff1a92ebe42ef7413a98b39b771b8759ad385ba3fb037898935cd09e}.
Its precedence and thresholds are explicit, but three execution details are not:

\begin{enumerate}
\item the score string uses \texttt{scoring\_staleness}, which has no JSON
definition;
\item same-precedence candidates have no deterministic tie-break; and
\item normal and META state transitions are not encoded in the policy.
\end{enumerate}

The current next choice remains unambiguous because area 13 is the unique
severe candidate.  To test the latent defect, the gate raises area 12 to the
same age.  The current artifact then returns a candidate set $\{12,13\}$ rather
than a total decision.  The proposal-only file
\texttt{POLICY\_PROPOSAL\_v1.1.json} resolves ties by age, current score, and
area ID, after defining the state transitions and scoring age.  It selects area
13 in the adversarial state and preserves the real Run 119/120 replay.

This does not justify adoption.  A different tie-break could be preferable,
and two-run no-drift evidence is much weaker than a full historical replay.
The proposal remains blocked on independent accountable review, schema
validation, an append-only machine transition ledger, and regression tests.

\section{Verification and negative controls}

The exact standard-library program \texttt{verification.py} was written and
executed before the paper.  Its complete output is
\texttt{verification\_output.txt}.  It passed 22 of 22 checks.

The gate verifies Eq.~\eqref{eq:gap} on a hand-auditable state, exhausts 78,100
combinations with two to six components, and samples 10,000 states with two to
32 components and a forced unique maximum.  It tests tied maxima, $\beta=0$,
stable log-sum-exp evaluation to $\beta=100$, the soft-maximum bounds, the live
META checksum, the unique next area, lens rotation, and the proposal replay.

The most important negative control assigns service costs 100 and 1 to ages 10
and 9.  After subtracting cost from one-step potential reduction, serving the
younger cheap component dominates.  Max-age is therefore not a general
cost-sensitive policy.  The control is not a failure of Theorem 1; it marks the
boundary of what Theorem 1 optimizes.

These finite checks reduce transcription, overflow, state, and implementation
risk.  They are not the formal proof.  The separate pinned Lean artifact
\texttt{SchedulerFreeEnergy.lean} proves the following five frozen targets:
\path{reset_partition_function_identity},
\path{max_age_minimizes_next_logsumexp},
\path{unique_max_age_unique_minimizer},
\path{tied_maxima_equal_next_potential}, and
\path{logsumexp_bounds_max_age}.  It also proves the concrete witness
\path{scheduler_free_energy_nonvacuous}.  An independent local audit rebuilt
the project in 1,895 jobs, confirmed statement freeze, zero sorries and escape
constructs, non-vacuity, and only \texttt{propext},
\texttt{Classical.choice}, and \texttt{Quot.sound}.  The formal result covers
the declared one-step equal-cost model; it does not prove the live policy-file
diagnosis, empirical productivity, or adoption authority.

\section{Evaluation plan}

A policy evaluation needs a machine-readable history, not the legacy narrative
alone.  Starting from a frozen state, an independent reference implementation
and the proposed selector should replay every directive.  Untied historical
decisions and resulting hashes must match exactly; tied cases must be surfaced,
not silently resolved.  Before/after evaluation should preregister maximum age,
tier-floor violations, attention diversity, inbound-flag debt, override rate,
and scientific-information proxies.  Reject the proposal if it changes an
untied decision, hides a prose dependency, increases maximum age, or reduces
diversity without a declared compensating benefit.

Scientific usefulness requires a harder model.  Add unequal service costs,
uncertain returns, delayed cross-pollination, and nonstationarity, then compare
max-age, score, risk-sensitive index, and curriculum policies on held-out
regimes.  The classical AoI, MaxWeight, Whittle, and learning-while-scheduling
literatures are the appropriate baselines, not internal META prose.

\section{Discussion and limitations}

The positive result is modest but useful: the severe-starvation action has an
exact declared objective under the implemented equal-cost transition.  Its
benefit does not depend on tuning $\beta$ when the maximum is unique.  The
negative result is operationally more important: a versioned policy can still
be incomplete if a symbol, state transition, or tie-break lives only in the
reader's assumptions.

The certificate does not establish optimal threshold 22, long-run stability,
scientific importance, novelty, external validity, or improved research output.
The log-sum-exp potential is selected for transparency, not inferred from data.
Tiers and flags may encode real value that age omits.  META scheduling also
consumes a slot without resetting an area, and future policy must evaluate that
cost rather than treat the cadence as physically necessary.

Nearest-work search found extensive prior theory for age, queue, index, and
curriculum scheduling.  No priority claim is made.  The internal fingerprint
index contains no exact match, but absence from that index cannot establish
novelty.  The paper may enter the standalone Viridis Research Operations
methods corpus.  The policy proposal itself remains on HOLD until full
historical replay, schema validation, accountable review, and separately
authorized adoption.

\section{Reproducibility and declarations}

\textbf{Reproducibility.} Run \texttt{python3 verification.py} with Python
3.11 or later; no third-party package or external data is required.  The seed,
tolerance, hypotheses, falsifier, source hashes, and full result are printed.
The manuscript compiles from \texttt{paper.tex} with Tectonic 0.17.0.

\textbf{Data and code availability.} All synthetic states, code, proposal
metadata, and captured outputs are in the immutable Run 119 package.  No
external scientific dataset was generated or analyzed.

\textbf{Funding.} No specific external funding supported this run.

\textbf{Conflicts of interest.} Justin D. Hart is founder of Viridis LLC,
which may develop research-governance or assurance tools.  The separate
business note reports zero signed revenue and no named customer validation.

\textbf{Accountable authorship.} Justin D. Hart is the accountable human
author and retains responsibility for scientific claims, corrections, and any
future policy or submission.  No AI system is represented as a journal author.

\textbf{Substantive AI-use disclosure.} Under Justin D. Hart's direction,
OpenAI Codex searched the local Viridis corpus and current primary literature,
preregistered the hypotheses and falsifiers, wrote and executed the gate,
audited the policy schema, drafted the manuscript, and assembled the package.
Historical Claude-authored Viridis artifacts were read for context with
attribution preserved.  Aristotle served as the automated theorem-proving
system for the six frozen Lean targets; Codex then rebuilt and audited the
returned artifact, checked statement preservation and scope, and reconciled
this final manuscript.  AI systems are research tools, not authors or policy
authorities.  No external peer review has occurred.

\textbf{Stage and route.} The package exits as
\texttt{FINAL\_POST\_ARISTOTLE\_RECONCILED}.  Release route: standalone
Viridis Research Operations methods corpus.  The policy proposal remains
\texttt{PROPOSAL\_ONLY\_NOT\_ADOPTED}; publication does not authorize its use.
No Intelligence Bound spine effect is proposed.

\begin{thebibliography}{9}
\bibitem{kaul2012}
S. Kaul, R. Yates, and M. Gruteser,
``Real-time status: How often should one update?''
\emph{IEEE INFOCOM}, 2731--2735 (2012).
\href{https://doi.org/10.1109/INFCOM.2012.6195689}{doi:10.1109/INFCOM.2012.6195689}.

\bibitem{kadota2018}
I. Kadota, A. Sinha, E. Uysal-Biyikoglu, R. Singh, and E. Modiano,
``Scheduling policies for minimizing age of information in broadcast wireless networks,''
\emph{IEEE/ACM Trans. Networking} \textbf{26}, 2637--2650 (2018).
\href{https://doi.org/10.1109/TNET.2018.2873606}{doi:10.1109/TNET.2018.2873606}.

\bibitem{tassiulas1992}
L. Tassiulas and A. Ephremides,
``Stability properties of constrained queueing systems and scheduling policies for maximum throughput in multihop radio networks,''
\emph{IEEE Trans. Automatic Control} \textbf{37}, 1936--1948 (1992).
\href{https://doi.org/10.1109/9.182479}{doi:10.1109/9.182479}.

\bibitem{yang2023}
Z. Yang, R. Srikant, and L. Ying,
``Learning while scheduling in multi-server systems with unknown statistics: MaxWeight with discounted UCB,''
\emph{AISTATS, PMLR} \textbf{206}, 4275--4312 (2023).
\href{https://proceedings.mlr.press/v206/yang23d.html}{PMLR record}.

\bibitem{cohen2024}
O. Cohen, R. Meir, and N. Weinberger,
``Statistical curriculum learning: An elimination algorithm achieving an oracle risk,''
\emph{COLT, PMLR} \textbf{247}, 1142--1199 (2024).
\href{https://proceedings.mlr.press/v247/cohen24a.html}{PMLR record}.
\end{thebibliography}

\end{document}
